\chapter{Joint Application--Network Optimization over Multipath}
\label{ch:multipath}

\Cref{ch:single_path} chose a single path for a flow and found its ceiling: when the shared bottleneck
saturates, even the best single path cannot carry the demand. This chapter starts from a more ambitious
premise. SCION offers not one path but several, so why commit to one? And if Mortise~\cite{Shen2026} can
tune a \emph{transport mechanism} to the application, why not tune the \emph{application itself} in the
same breath? We design and evaluate a system that (i) uses multiple paths at once via Multipath
QUIC~\cite{DeConinck2017} and (ii) jointly optimizes what the application sends -- the video encoder's
target bitrate -- and how the network spreads it -- the per-path traffic split -- to maximize the quality
of experience of a real-time video call. \Cref{sec:mp:design} presents the design, \cref{sec:mp:impl} the
testbed, and \cref{sec:mp:eval} the evaluation, whose central finding is that neither decision can be
made well alone, and that the value of making them jointly is governed by how unevenly capacity is
distributed across the paths and by how strictly the application prices delay.

\gap{this reframing has been propagated to the abstract, \cref{sec:disc:interpretation},
\cref{sec:analysis:generalization}, and \cref{ch:conclusion}, but \textbf{one site remains and it needs
an authorial decision rather than a rewording}: contribution~5 of \cref{sec:intro:contributions} is
titled ``characterization of the volatility driver for multipath control'', which \emph{names} the
claim. On the evidence of \cref{sec:p2eval:regimes} volatility is not the driver --- it is one mechanism
that makes path headroom uneven, and a stationary asymmetric topology produces the same collapse. Either
rename it (``characterization of when joint multipath control pays off'') or keep the volatility
framing and weaken \cref{sec:p2eval:regimes} to match; the two currently disagree, and the contributions
list is the more visible of the two.}

\gap{\textbf{Chapter-level, highest priority: this chapter currently has no intent in it.}
\Cref{sec:p1design:conditioning} closes by promising that ``\cref{ch:multipath} pushes the same
principle one step further, letting intent shape not only the network's choice but the application's own
output'', and \cref{sec:rel:gcrl} already asserts the cancellation result for a comparison-based
``argmax \emph{or softmax}'' selector. Neither is delivered here: the reward weights below are fixed
constants and no intent vector enters any observation. The implementation does have an intent mechanism
-- \texttt{rl-mpquic/configs/profiles/\{interactive,presenter,passive\}.yaml}, each bundling reward
weights \emph{and} a \texttt{deadline\_ms} of 180/400/800 -- but it is realized as three
\emph{separately trained} policies, which is precisely the per-objective model zoo
\cref{ch:single_path} argues against.

Two ways out, and the chapter must take one by \textbf{2026-08-08} --- late enough for the conditioning
work to have reported, early enough to rewrite the chapter's framing before the experiment freeze on
2026-08-18.

\emph{(i) Lead with intent conditioning.} Implement the role-conditioned single policy
(\texttt{rl-mpquic/docs/ROADMAP.md} 4.3, specified as work item W1 in
\texttt{rl-mpquic/docs/THESIS\_BRIEF.md}) and make it this chapter's centrepiece, together with the
softmax proposition below. This is what makes the chapter the core of the thesis, and it is what the
title, \cref{sec:p1design:conditioning}, and \cref{sec:intro:contributions} currently promise.

\emph{(ii) Lead with joint co-optimization.} If the conditioned policy does not clear its bar, the
chapter's contribution is the \emph{joint} framing rather than the conditioned one, and that claim is
already evidenced: the two-sided ablation of \cref{sec:p2eval:ablation} shows both halves losing on both
backends under every intent, and no prior system co-optimizes rate and split --- adaptive-bitrate work
adapts rate over one exogenous path, multipath schedulers schedule an exogenous load
(\cref{sec:rel:transport,sec:rel:abr}). Note this is what the printed subtitle already promises. Under
(ii) the multipath system is presented as serving three application intents with per-intent
specialists, the bridging sentence in \cref{sec:p1design:conditioning} is softened rather than deleted,
and the conditioning attempt is reported as a finding: an objective can price a term too weakly to
generate the gradient that would train against it, which is the mechanism already visible in the failed
\texttt{passive} diagonal of \cref{tab:p2transfer}.

Under \emph{both} options the softmax proposition ships, since it is a mathematical result rather than
an experimental one, and the evaluation backbone is unchanged --- so this decision governs the
chapter's framing, not which measurements are made.}

\section{System Design: Co-Optimizing Rate and Traffic Split}
\label{sec:mp:design}

\Cref{sec:p2design:motivation} motivates the joint, multipath framing and the strawmen it improves on.
\Cref{sec:p2design:hierarchy} presents the two-timescale hierarchical agent pair.
\Cref{sec:p2design:mdp} specifies the observations, actions, and rewards.
\Cref{sec:p2design:coupling} explains the coupling that keeps the two agents cooperating rather than
working at cross purposes. \Cref{sec:p2design:scoring} extends the scoring idea of \cref{ch:single_path}
from a discrete choice to a continuous split over a variable, possibly changing, path set.

\subsection{Motivation and Strawmen}
\label{sec:p2design:motivation}

Interactive video is a demanding, ubiquitous workload with no playback buffer to hide variance: a frame
that misses its deadline is simply lost. Its QoE depends on two decisions that are usually made in
isolation. The \emph{application} chooses an encoding bitrate -- too high and the network drops frames,
too low and quality suffers. The \emph{network} chooses how to carry the resulting bytes -- over one path
or spread across several. Optimizing either alone is a strawman:

\begin{description}
  \item[Application-only.] Adapt bitrate to the aggregate network but carry it over a fixed or naive
  path assignment. This wastes SCION's path diversity and, as our results show, collapses when the paths
  it depends on degrade.
  \item[Network-only.] Schedule cleverly across paths but treat the offered load as exogenous. The
  scheduler cannot relieve congestion it is itself creating by admitting too high a bitrate.
\end{description}

The two decisions are coupled through the same bottleneck: bitrate sets the load the scheduler must
place, and the scheduler's success sets the effective capacity the bitrate controller sees. Optimizing
them jointly is the design's premise. The challenge is that they naturally operate on \emph{different
timescales} -- an encoder should not change bitrate every frame, but a scheduler must react per frame --
which motivates a hierarchy rather than a single monolithic agent.

\subsection{Two-Timescale Hierarchical Agents}
\label{sec:p2design:hierarchy}

We split control into two cooperating agents that act at different cadences, a ``brain and body''
decomposition.

\begin{description}
  \item[Application agent (slow).] Every \SI{1}{\second} (once per 30 frames at 30\,fps) it observes
  an aggregate summary of the connection and outputs a single target encoder bitrate in the range
  \SIrange{300}{6000}{\kilo\bit\per\second}. It is deliberately \emph{horizon-agnostic} -- it never
  observes how far along the call is -- so the same policy applies to a call of any length, an essential
  property for a controller meant to run in unbounded real sessions.
  \item[Transport agent (fast).] Every frame (${\approx}\SI{33}{\milli\second}$) it observes per-path
  transport state \emph{and the application agent's current target bitrate}, and outputs a split of that
  frame's bytes across the available paths.
\end{description}

Both are trained with soft actor-critic (\cref{sec:background:rl}), whose continuous action support and
sample efficiency suit these controls and the cost of simulator rollouts.

\subsection{Observations, Actions, and Rewards}
\label{sec:p2design:mdp}

\paragraph{Application observation and action.} The application agent sees a small fixed-size vector:
normalized current bitrate, smoothed RTT, jitter, loss, and throughput. Its action is one continuous
scalar in $[-1,1]$, mapped affinely to a target bitrate in \SIrange{300}{6000}{\kilo\bit\per\second}.

\paragraph{Transport observation and action.} The transport agent sees a global block (target bitrate,
RTT, loss, throughput) and, for each path, a feature vector (congestion window, smoothed RTT, buffer
occupancy, per-path throughput, per-path loss). Its action is an $N$-vector passed through a softmax to
yield a split $(w_1,\dots,w_N)$ that is non-negative and sums to one -- a valid allocation by
construction.

\paragraph{Rewards.} The rewards encode the QoE objective and the division of responsibility between the
agents. The application's per-window QoE reward composes a quality term with penalties:
\begin{equation}
\mathrm{QoE} = a\cdot\frac{\mathrm{VMAF}(\cdot)}{100}
             - b\cdot\widetilde{\mathrm{lat}}
             - c\cdot\widetilde{\mathrm{jit}}
             - d\cdot\mathrm{loss}
             + e\cdot(1-\mathrm{loss})\cdot\widetilde{\mathrm{util}},
\qquad \text{clipped to } [-2, 1],
\label{eq:p2qoe}
\end{equation}
where $\widetilde{\mathrm{lat}}$ and $\widetilde{\mathrm{jit}}$ are latency and jitter normalized by
$\mathit{lat}_{\mathrm{norm}}$ and $\mathit{jit}_{\mathrm{norm}}$ and clipped, and
$\widetilde{\mathrm{util}}$ is the achieved bitrate normalized by
\SI{3300}{\kilo\bit\per\second} -- the knee of the quality curve. The utilization term restores the
rate-control gradient under a quality model that saturates in bitrate; it carries $e = 0.25$ in the
dynamic scenario and is harmless under the default logarithmic curve. The quality term uses a VMAF
surrogate (\cref{sec:mp:impl}); when a learned surrogate that already accounts for loss is used, the
explicit loss penalty is dropped to avoid double-counting.

The weight vector $(a,b,c,d,e)$ together with the frame deadline is not a constant of the system but
the \emph{application intent} this thesis is about, and the three profiles evaluated in
\cref{sec:mp:eval} instantiate it at three points of the ITU-T G.1010 use-case space
(\cref{tab:p2profiles}). The transport agent receives a per-frame reward that credits delivery without
quality, since it does not control bitrate:
\begin{equation}
R_{\text{transport}} = (1 - \mathrm{loss}) - b\cdot\widetilde{\mathrm{lat}} - c\cdot\widetilde{\mathrm{jit}},
\qquad \text{clipped to } [-2, 1].
\label{eq:p2transport}
\end{equation}

\begin{table}[t]
  \centering
  \caption{The three application intents evaluated in \cref{sec:mp:eval}, each a bundle of QoE weights
  (\cref{eq:p2qoe}), normalizers, and a frame deadline, keyed to ITU-T G.1010 use-case classes. Source:
  \texttt{rl-mpquic/configs/profiles/*.yaml}. The delay weight $b$ falls and the deadline widens as the
  role relaxes; loss stays dominant for the two content-critical roles. These play the part the reward
  profiles of \cref{tab:app:profiles} play in \cref{ch:single_path}.}
  \label{tab:p2profiles}
  \begin{tabular}{llccccc}
    \toprule
    Intent & G.1010 class & Deadline & $b$ & $c$ & $d$ & $\mathit{lat}_{\mathrm{norm}}$ \\
    \midrule
    \texttt{interactive} & conversational video & \SI{180}{\milli\second} & 0.25 & 0.25 & 2.0 & \SI{200}{\milli\second} \\
    \texttt{presenter}   & one-way video        & \SI{400}{\milli\second} & 0.20 & 0.20 & 2.0 & \SI{400}{\milli\second} \\
    \texttt{passive}     & background/streaming & \SI{800}{\milli\second} & 0.10 & 0.30 & 0.7 & \SI{800}{\milli\second} \\
    \bottomrule
  \end{tabular}
\end{table}

\gap{document the provenance of these weights, which is a methodological contribution the chapter
currently omits entirely. The delay anchors come from ITU-T G.114 (\SI{150}{\milli\second} preferred,
\SI{400}{\milli\second} unacceptable) and the jitter normalizer from Y.1541 class-0 IPDV; the pair
$(b,d)$ is then \emph{fit} rather than hand-set, by grid-correlating the linear reward of
\cref{eq:p2qoe} against an ITU-T G.1070 opinion-model oracle over a trace corpus
(\texttt{rl-mpquic/src/ns3env/g1070.py}, \texttt{scripts/calibrate\_reward.py}; protocol in
\texttt{docs/REWARD\_TUNING.md} \S5). Report two things with it: that the fit is corpus-insensitive
(synthetic and real-trace corpora both land on $b = c = 0.25$, $d = 2.0$), and the \emph{stability gate}
of \S5C -- a weight set that changes which policy is optimal is rejected regardless of its correlation.
That gate is what the static scenario fails (\cref{sec:p2eval:results}), and stating it before the
results turns an awkward negative into a methodological point.}

\subsection{Coupling the Two Agents}
\label{sec:p2design:coupling}

The design's subtlety is keeping the agents cooperative. Two mechanisms couple them.

\paragraph{Observation coupling.} The transport agent observes the application's current target bitrate.
It therefore schedules \emph{knowing} how much load it must place, and can, for instance, spread an
aggressive bitrate across more paths rather than overrunning one. This is what makes the pair hierarchical
rather than two agents optimizing in ignorance of each other.

\paragraph{Reward decomposition.} The application owns quality (the VMAF term) and is credited for the
latency, jitter, and loss experienced over the window its bitrate governed; the transport agent is
credited only for how well it delivered each frame. This assigns each decision the consequences it can
actually control -- the encoder is not blamed for a poor split, nor the scheduler for an over-ambitious
bitrate -- while the shared bottleneck still couples them: a bad split shows up as latency/loss in the
application's reward, and an over-high bitrate shows up as loss the transport agent cannot schedule away.

\subsection{Continuous Scoring over a Variable Path Set}
\label{sec:p2design:scoring}

The transport agent inherits the single-path selector's problem -- a variable, possibly changing number
of paths -- and its solution, adapted from a discrete choice to a continuous allocation. A shared per-path
encoder maps each $(\text{global}, \text{path}_i)$ pair to a logit; the logits are softmaxed into the
split, so the policy is permutation-equivariant and independent of the path count. To handle paths that
appear and disappear, each path carries a \emph{liveness} flag and the state includes a binary mask; the
softmax excludes dead paths, so traffic is never allocated to a path that is gone. The critic aggregates
per-path embeddings with a masked mean-pool -- a Deep Sets~\cite{Zaheer2017} operation -- yielding a value
estimate invariant to path ordering and count. This continuous scoring agent is the natural sibling of the
discrete scoring selector of \cref{ch:single_path}: the same architectural idea, once for ``pick a path''
and once for ``divide across paths.''

\gap{\textbf{The chapter's missing theoretical result, and the cheapest way to make it the core.}
\Cref{prop:p1cancel} shows that conditioning entering only path-independent terms cancels under the
$\arg\max$. The allocation here is a \emph{softmax} over per-path logits, and softmax is exactly
invariant to a constant added to every logit: $\operatorname{softmax}(z + c\mathbf{1}) =
\operatorname{softmax}(z)$ for any $c$. So the same result holds in the continuous case, and holds
\emph{more} sharply -- state it as a proposition alongside \cref{prop:p1cancel}. Three reasons it is
worth the page. First, it is strictly stronger than the discrete statement: the split is not merely
``not re-ranked'' but numerically identical, with no tie-breaking caveat and none of the trust-feature
leakage that softens \cref{prop:p1cancel} in \cref{sec:p1design:concat}. Second, it is verifiable to
machine precision rather than behaviorally: hold the path rows fixed, sweep the intent through the
global block alone, and the split must be identical to floating-point tolerance; route the intent into
the per-path rows instead and it must move. Third, it discharges a claim the thesis already makes on
credit -- \cref{sec:rel:gcrl} asserts the result for ``argmax \emph{or} softmax'' selectors and
nothing in the document currently proves the second half. Pair the proposition with the two-arm
experiment (global-only conditioning vs.\ per-path conditioning) and report behavioral divergence with
the metric of \cref{eq:p1divergence}, adapted from a discrete choice count to a distance between
splits. \texttt{rl-mpquic/docs/THESIS\_BRIEF.md} specifies this as work item W2.}

\section{Implementation: Multipath Real-Time Video Testbed}
\label{sec:mp:impl}

Realizing the design requires a testbed that is faithful enough to expose real multipath and
congestion behavior yet fast enough for reinforcement-learning rollouts. We resolve this tension with a
single data-plane abstraction served by two interchangeable backends (\cref{sec:p2impl:dataplane}): a
packet-level NS-3 simulator and a fast analytical mock. \Cref{sec:p2impl:rt} describes the real-time
video and multipath model, \cref{sec:p2impl:agents} the agent implementation and the VMAF surrogate, and
\cref{sec:p2impl:gaps} the known gaps that scope the current results.

\subsection{Data-Plane Abstraction and Two Backends}
\label{sec:p2impl:dataplane}

Both backends implement one abstract \texttt{DataPlane} interface, so the RL code, observations, and
rewards are written once and run unchanged against either. This lets us iterate quickly on the mock and
validate on the higher-fidelity simulator without divergence.

\begin{description}
  \item[NS-3 backend.] A packet-level simulation in which the connection's several paths are modeled as
  independent subflows over a topology of links, each with its own capacity, propagation delay, and
  active-queue management (FqCoDel). Realistic contention comes from OnOff UDP cross-traffic with
  exponentially distributed on/off periods, phase-shifted per path so the paths do not degrade in
  lockstep. Each video frame's bytes are striped across the subflows according to the transport agent's
  split, and delivery is tracked by a byte-watermark scheme -- avoiding per-packet tagging -- from which
  per-frame latency, jitter, and deadline-miss loss are measured. Python and NS-3 communicate through an
  ns3-ai shared-memory bridge.
  \item[Mock backend.] A pure-Python, per-seed-deterministic model in which each path has a capacity
  envelope (a smooth sinusoidal component plus bursty cross-traffic and noise) and a standing-queue model
  that produces bufferbloat when a path is overdriven. It runs orders of magnitude faster than NS-3 and is
  used for rapid training.
\end{description}

\paragraph{Non-stationary dynamics.} Both backends inject four mechanisms that repeatedly unbalance the
capacity available on each path: \emph{path churn} (paths appear and disappear, and bytes sent to a dead
path are lost), \emph{regime shifts} (abrupt changes in the best path's capacity), \emph{congestion bursts}
(transient per-path capacity collapses), and \emph{correlated failures} (shared-bottleneck groups that
degrade together). The NS-3 body implements churn with drop-all receive error models rather than by
collapsing the link rate, which would strand a packet mid-serialization, and tears down and reconnects
the affected subflow; regime, burst, and correlation scale the bottleneck rate and the cross-traffic
rate together. Behavioral parity between the two bodies is checked by a fixed-policy loss-regime guard
(\texttt{rl-mpquic/scripts/parity\_check.py}).

\gap{report the parity check as a result rather than leaving it as a build step, and state the two
places the bodies are deliberately not identical: they draw from independent random streams, so they are
behaviorally equivalent rather than frame-identical, and the mock's sender queue is unbounded where
NS-3's is not. The second asymmetry is not cosmetic -- it is why the mock's failure model is optimistic,
and it is the mechanism behind the \texttt{single} baseline inverting between the two bodies
(\cref{tab:p2ns3}). \texttt{rl-mpquic/docs/TUNING\_DYNAMICS.md} \S4 lists the structural differences.}

\subsection{Real-Time Video and Multipath Model}
\label{sec:p2impl:rt}

The video source mirrors a conferencing encoder: a fixed frame rate (30\,fps by default), variable
frame sizes (with periodic larger I-frames and per-frame jitter), no playback buffer, and a deadline
after which a late frame is discarded as lost. The aggregate target bitrate ranges over
\SIrange{300}{6000}{\kilo\bit\per\second}. The canonical topology is deliberately \emph{non-dominant}: no
single path can carry full-quality video, and the aggregate usable capacity sits below the quality knee,
so the system must genuinely combine paths and trade rate against delivery. For example, a four-path
configuration mixes paths of differing capacity and RTT -- including a higher-capacity but high-latency
``trap'' path -- under heavy, path-specific cross-traffic, so that neither the widest nor the
lowest-latency path is right on its own.

\subsection{Agents and QoE Surrogate}
\label{sec:p2impl:agents}

\paragraph{Agents.} Both agents are soft actor-critic learners with twin critics and automatic
temperature tuning. The transport agent has two variants: a \emph{flat} version whose actor and critic
consume a fixed-length concatenation of the global and per-path features (a fixed path count), and a
\emph{scoring} version implementing the permutation-equivariant, mask-aware design of
\cref{sec:p2design:scoring} with a shared per-path actor and a Deep Sets masked-mean-pool critic. The
hierarchical training loop runs the two agents at their respective cadences, storing each agent's
transitions in its own replay buffer and crediting the application agent over the window of frames its
bitrate governed.

\paragraph{VMAF surrogate.} The quality term of \cref{eq:p2qoe} needs a mapping from an encoding (and
network conditions) to perceptual quality. Two are supported: a default bitrate-only logarithmic curve fit
through published quality anchors, and a learned surrogate that interpolates a
quality-of-service$\rightarrow$VMAF table over bitrate, loss, delay, and jitter (produced by a sibling
data-generation project). When the learned surrogate is active, the explicit loss penalty in
\cref{eq:p2qoe} is dropped to avoid double-counting the loss it already reflects.

\subsection{Known Gaps}
\label{sec:p2impl:gaps}

The current testbed has four limitations that scope the results of \cref{sec:mp:eval} and define concrete
follow-up work.

\begin{itemize}
  \item \textbf{Transport realism.} Subflows are modeled with TCP-like congestion control rather than true
  Multipath QUIC; this preserves congestion, queueing, and loss dynamics but masks some network loss as
  latency, and forgoes QUIC's stream multiplexing, connection migration, and unreliable datagrams.
  Swapping in a real QUIC module is future work.
  \item \textbf{Surrogate coverage.} The learned VMAF surrogate's loss, delay, and jitter axes are
  currently under-sampled, so the quality term is effectively bitrate-dominated; a richer surrogate would
  activate those axes and let quality respond to delivery conditions directly. Every result in
  \cref{sec:mp:eval} therefore uses the bitrate-only logarithmic curve.
  \item \textbf{Training body.} All reported policies are trained on the mock backend and evaluated on
  both. That a mock-trained policy transfers is measured (\cref{sec:p2eval:ns3}); that the system
  \emph{learns} on the packet-level body is not yet.
  \item \textbf{Episode independence.} Episodes reset in-band to keep the shared-memory bridge alive,
  which leaves congestion state warm across resets, so episodes are not independent draws.
\end{itemize}

\gap{one further scoping item belongs here once measured: the application agent is horizon-agnostic by
construction (\cref{sec:p2design:hierarchy}), but every episode is \SI{30}{\second} long, so the
property is a design claim rather than an evaluated one. A single long-call rollout -- several minutes,
policy unchanged -- would convert it into a measurement, and it costs one evaluation run.}

\section{Experimental Evaluation: Regimes of Value}
\label{sec:mp:eval}

This section evaluates the multipath system and answers five questions in turn.
\Cref{sec:p2eval:setup} fixes the setup, the eight comparison points, and the statistical protocol.
\Cref{sec:p2eval:results} asks whether the learned pair beats every baseline, and under which
application intents. \Cref{sec:p2eval:regimes} asks the question the chapter is named for --
\emph{when} does learned scheduling earn its keep -- and answers it with a mechanism rather than a
scenario label. \Cref{sec:p2eval:transfer} asks whether the three intents produce genuinely different
policies or merely relabel the reward. \Cref{sec:p2eval:ns3} asks whether any of it survives on a
packet-level transport stack, and \cref{sec:p2eval:ablation} isolates the individual design choices,
including the ones that did not pay off.

\subsection{Setup, Baselines, and Statistical Protocol}
\label{sec:p2eval:setup}

\paragraph{Scenarios.} The \emph{dynamic} scenario is the headline test: six candidate paths, none
individually sufficient, with a latency trap, a shared-bottleneck pair, and all four non-stationary
mechanisms of \cref{sec:p2impl:dataplane} enabled, so that the identity and quality of the best paths
change during a call. The \emph{static} scenario is a three-path stationary topology retained as a
contrast case; \cref{sec:p2eval:regimes} explains why it is the less informative of the two.

\paragraph{Metric.} The primary metric is mean per-window QoE (\cref{eq:p2qoe}) under the intent being
evaluated; we also report the underlying achieved bitrate, latency, loss, and VMAF. Because each intent
carries its own reward, QoE is comparable \emph{within} an intent and never across intents --- all
comparisons below are read down a column.

\paragraph{Baselines and ablations.} Eight policies are compared. Four are fixed scheduling rules
driving a common reactive bitrate controller, so that the comparison isolates \emph{scheduling}:
\emph{even} split over live paths, \emph{single} highest-throughput path, \emph{proportional} allocation
by recent per-path throughput, and \emph{random} (a fresh Dirichlet draw per frame). Under dynamics all
four are mask-aware, operating only over active paths, so they are not handicapped by churn they cannot
see. A fifth, \emph{WebRTC}, is the strongest external comparison point: a stateful Google Congestion
Control bitrate estimator driving a proportional split, i.e.\ what a deployed real-time video stack does
today. Two ablations decompose the learned pair itself: \emph{app-only} keeps the learned bitrate
controller and replaces the transport agent with an even split, and \emph{path-only} keeps the learned
split and replaces the application agent with the same GCC estimator. Together the two ablations answer
which half of the hierarchy carries the result.

\paragraph{Statistical protocol.} Every comparison is graded over a grid of training checkpoints
$\times$ evaluation seeds, five episodes per cell, with all policies seeing identical conditions within
a cell. Three quantities accompany every claim: the 95\% confidence interval over evaluation seeds, the
\emph{training-seed spread} $\mathit{SD}_{\mathrm{train}}$ (the standard deviation of per-checkpoint
mean QoE), and a per-cell ranking-stability check that reports whether the learned pair is first in
\emph{every} cell rather than merely on average. The middle quantity is the one that matters, and it
deserves stating as a methodological position rather than a footnote: a margin smaller than
$\mathit{SD}_{\mathrm{train}}$ is a statement about which training seed was drawn, not about the method.
Measuring it overturned three of four verdicts we had previously drawn from single-checkpoint
evaluations, and every such reversal is reported below.

\subsection{Does the Learned Pair Win, and Under Which Intent?}
\label{sec:p2eval:results}

\Cref{tab:p2results} reports the dynamic scenario under each of the three application intents of
\cref{tab:p2profiles}.

\begin{table}[t]
  \centering
  \caption{Mean QoE (\cref{eq:p2qoe}) on the dynamic six-path scenario, by policy and application
  intent, over three evaluation seeds $\times$ five episodes. Each column is evaluated under its own
  reward and deadline, so columns are not comparable to one another. The learned pair is first in every
  column and in every individual cell. Source: \texttt{rl-mpquic/runs/ab-prof-*/ab\_summary.csv}.}
  \label{tab:p2results}
  \begin{tabular}{lccc}
    \toprule
    Policy & \texttt{interactive} & \texttt{presenter} & \texttt{passive} \\
    \midrule
    Learned pair              & \textbf{0.636} & \textbf{0.736} & \textbf{0.808} \\
    Path-only (GCC rate)      & 0.575 & 0.615 & 0.671 \\
    WebRTC (GCC + proportional) & 0.570 & 0.601 & 0.667 \\
    Single best-active path   & 0.456 & 0.496 & 0.512 \\
    App-only (learned rate, even split) & 0.247 & 0.407 & 0.657 \\
    Proportional              & 0.238 & 0.275 & 0.332 \\
    Even split                & 0.126 & 0.158 & 0.223 \\
    Random split              & 0.118 & 0.154 & 0.217 \\
    \bottomrule
  \end{tabular}
\end{table}

The learned pair leads every column, and the ranking is stable: it is first in each of the nine
evaluation cells, not merely on the mean. Its margin over the strongest comparison point --- the
path-only ablation, which is itself a learned scheduler --- is 0.061 (10.7\%) under
\texttt{interactive}, 0.121 (19.6\%) under \texttt{presenter}, and 0.137 (20.4\%) under
\texttt{passive}. The deployed-stack comparison is the more legible one: the learned pair beats a
WebRTC-style GCC controller with a proportional split by 11.6\% to 22.5\% depending on the intent.

The policies also \emph{behave} differently under different intents, which is what the intent mechanism
is for. Delivered latency rises from \SI{49.0}{\milli\second} (\texttt{interactive}) to
\SI{57.6}{\milli\second} (\texttt{presenter}) to \SI{93.1}{\milli\second} (\texttt{passive}), and the
achieved bitrate rises with it from 2626 to 3304 to \SI{3400}{\kilo\bit\per\second}, lifting VMAF from
79.5 to 85.3 to 86.1. As the intent relaxes its delay budget, the policy spends the slack on quality.
This is the multipath analogue of the intent alignment demonstrated in \cref{sec:p1eval:intent}, and it
is the closest the chapter currently comes to the thesis's central claim.

\gap{this table rests on \emph{one training seed per intent}, and the \texttt{interactive} checkpoint
was trained for 200 episodes against the other two at 100 --- so the cross-intent comparisons of
convergence are confounded, even though the within-column baseline comparisons are not. The matched
budget rerun (three intents $\times$ three seeds at 100 episodes) is
\texttt{rl-mpquic/docs/RUNBOOK.md} Step~1 and was executing at the time of writing
(\texttt{runs/prof2-*}). Replace every number in \cref{tab:p2results} from
\texttt{runs/ab-prof2-*/ab\_summary.csv} when it completes, and add the $\mathit{SD}_{\mathrm{train}}$
column: the 10--20\% margins here are comfortably outside the 0.013--0.151 range measured elsewhere,
but they should be shown to be, not assumed to be.}

\gap{no figure in this chapter yet, and this is the section that most needs two. (i) A within-episode
trace: the per-path split alongside path liveness and per-path capacity, showing the transport agent
reallocating away from a path as it churns or bursts --- this is the mechanism claim of the whole
chapter and it is currently asserted in prose. (ii) A grouped bar or slope chart of
\cref{tab:p2results}, which reads poorly as an eight-row table. Both render from an existing
\texttt{evaluation\_results.json} through \texttt{rl-mpquic/evaluation/generate\_figures.py}
(\texttt{figure6\_timeseries}, \texttt{figure7\_split\_behavior}, \texttt{figure1\_qoe}); note that no
figure has yet been generated for \emph{any} of the current results, because \texttt{-\/-figures} was
never passed.}

\subsection{When Does Splitting Help? The Governing Quantity}
\label{sec:p2eval:regimes}

The chapter's title promises regimes of value, and the honest form of that claim is not ``learned
scheduling helps when the network is dynamic.'' It is sharper, and it is mechanical. The quantity that
governs whether a scheduler has anything to do is the relationship between the \emph{share an even
split places on each path} and the \emph{residual capacity of the weakest path}. Where every path can
absorb its even share, any allocation works and the learned agent has nothing to earn. Where some path
cannot, the split decides whether the frame arrives. Volatility matters because churn, regime shifts,
and bursts are mechanisms that \emph{destroy} that headroom homogeneity transiently --- but they are not
the only ones, and naming them as the axis mistakes a cause for the variable.

Two observations support the reframing, and both cut against the simpler story.

\paragraph{A static network is not a safe network.} On the static three-path scenario, the app-only
ablation does not merely lose --- it collapses to $-0.396$ QoE at \SI{696}{\milli\second} latency and
52\% loss, because the topology is \emph{asymmetric} (8/4/\SI{2}{\mega\bit\per\second}) and an even
split overdrives the weakest path continuously. Staticness bought it nothing. Conversely, a static
topology of four \emph{near-equal} paths puts app-only within 2\% of the full system. It is the
homogeneity of headroom, not the stationarity of the network, that decides.

\paragraph{The static scenario cannot discriminate at all.} On that same three-path scenario the
learned pair scores 0.700 against 0.691 for the single-best-path baseline --- a margin of 0.009 against
an $\mathit{SD}_{\mathrm{train}}$ of 0.017. The learned pair is first in six of nine cells and loses all
three cells of one training seed. The correct reading is not that scheduling is unnecessary when paths
are stable; it is that this scenario \emph{cannot tell}, because the margin is smaller than the spread
between training runs of the same method. A practitioner benchmarking a multipath scheduler only on a
stable topology would conclude it does not work, and would be measuring their own seed draw.

\paragraph{Intent strictness is the second axis.} The app-only ablation recovers 39\%, 55\%, and 81\%
of the full system's QoE under \texttt{interactive}, \texttt{presenter}, and \texttt{passive}
respectively (\cref{tab:p2results}). Holding the network fixed and varying only what the application
asks for, the scheduler's contribution more than doubles as the delay budget tightens. This is a
mechanism-level result rather than an average: the stricter the intent, the less slack there is to
absorb a bad allocation, so the more the allocation matters. It is also the point at which this chapter
rejoins \cref{ch:single_path} --- the value of the control machinery is itself a function of the intent.

\gap{turn the two paragraphs above into a measured two-dimensional map, which would be this chapter's
counterpart to \cref{fig:p1eval:ceiling} and its strongest single figure. Rows: headroom heterogeneity,
swept by scaling the per-path capacity spread and the dynamics hazard rates
(\texttt{rl-mpquic/configs/dynamic.yaml}, \texttt{dynamics:} block, and the held-out
\texttt{dynamic\_ood\_\{harsh,mild\}.yaml}). Columns: the three intents of \cref{tab:p2profiles}. Cell
value: the app-only recovery ratio, which is the cleanest scalar summary of ``how much did the
scheduler contribute here''. Nine reference points already exist in
\texttt{rl-mpquic/docs/TUNING\_DYNAMICS.md} \S4 and can seed the grid; the full sweep is evaluation-only
against checkpoints already on disk. Specified as work item W4 in
\texttt{rl-mpquic/docs/THESIS\_BRIEF.md}.}

\subsection{Do the Intents Produce Different Policies?}
\label{sec:p2eval:transfer}

\Cref{tab:p2results} shows that policies trained under different intents behave differently, but that is
also consistent with the three intents simply being easier or harder to satisfy. The discriminating test
is to evaluate every intent-trained policy under every intent's reward, and read down each column: if
the intents produce genuinely specialized policies, each column should be won by its own diagonal
entry. \Cref{tab:p2transfer} reports it. This is the multipath counterpart of the intent-alignment
matrix of \cref{fig:p1eval:heatmap}, and unlike that one it does not come out clean.

\begin{table}[t]
  \centering
  \caption{Intent-transfer matrix: mean QoE of each intent-trained policy (rows) evaluated under each
  intent's own reward and deadline (columns), on the dynamic scenario. Read down a column. The diagonal
  wins for \texttt{interactive} and \texttt{presenter}, and \emph{loses} for \texttt{passive}. Source:
  \texttt{rl-mpquic/runs/xrole-eval-*/trainedrole/seed-*-t*/evaluation\_results.json}.}
  \label{tab:p2transfer}
  \begin{tabular}{lccc}
    \toprule
    & \multicolumn{3}{c}{Evaluated under} \\
    \cmidrule(l){2-4}
    Trained under & \texttt{interactive} & \texttt{presenter} & \texttt{passive} \\
    \midrule
    \texttt{interactive} & \textbf{0.636} & 0.686 & 0.767 \\
    \texttt{presenter}   & 0.628 & \textbf{0.736} & \underline{0.828} \\
    \texttt{passive}     & 0.613 & 0.723 & 0.808 \\
    \bottomrule
  \end{tabular}
\end{table}

Two of the three diagonals hold, and the specialization is real rather than a relabeling of the reward:
the \texttt{presenter} policy beats both others under the presenter reward by 0.013 and 0.050 on every
seed, and loses to the \texttt{interactive} policy under the interactive reward. A policy trained for an
intent is, in those two cases, the right policy for it.

The third refutes the claim, and the failure is informative. Under \texttt{passive}'s own reward the
\texttt{passive}-trained policy is beaten by the \texttt{presenter}-trained one by 0.020 on every seed,
while the presenter policy simultaneously delivers \emph{lower} latency (\SI{65}{\milli\second} against
\SI{93}{\milli\second}) and \emph{higher} bitrate (3575 against \SI{3400}{\kilo\bit\per\second}): it is
better on passive's own terms by every measure. The mechanism is visible in the intent itself.
\texttt{passive} prices delay at $b = 0.1$, which is small enough that a lazy scheduler is barely
penalized during training, so learning settles in a slack local optimum that a policy trained under
pressure simply dominates. The lesson generalizes beyond this system: \emph{a relaxed objective is not
an easy one to learn}, and a weight can be too small to generate the gradient that would train against
it.

\gap{this matrix is the strongest available argument for the conditioned policy of
\cref{sec:p2design:scoring}, and the chapter should make that argument explicitly. If a single
intent-conditioned policy is trained, the bar it must clear is not the diagonal but the \emph{best entry
in each column} --- 0.636, 0.736, and 0.828 --- where the third is held by an off-diagonal policy. That
also means the conditioned policy has a plausible free win in the \texttt{passive} column, since the
specialist there is known to be beatable. Report the comparison as a fourth row of
\cref{tab:p2transfer}.}

\gap{re-derive this matrix from the matched-budget checkpoints (\texttt{runs/xrole2-s\{1,2,3\}-eval-*},
\texttt{RUNBOOK.md} Step~1) before treating it as settled. Two confounds currently limit it: the
\texttt{interactive} row had twice the training budget of the others, and every cell is a single
training seed, so the $+0.008$ interactive diagonal margin sits inside the measured
$\mathit{SD}_{\mathrm{train}}$ range. The \texttt{passive} result is the robust one --- it reproduces on
3 of 3 evaluation seeds with a tight interval --- but whether it reproduces across \emph{training} seeds
is exactly what the rerun answers.}

\subsection{Does It Survive a Packet-Level Transport Stack?}
\label{sec:p2eval:ns3}

Every result so far is measured on the analytical mock. \Cref{tab:p2ns3} evaluates a policy trained
entirely on the mock, without retraining or fine-tuning, on the NS-3 body --- real TCP subflows, FqCoDel
queues, and packet-level cross traffic.

\begin{table}[t]
  \centering
  \caption{Zero-shot transfer to the packet-level NS-3 backend: a mock-trained policy evaluated on the
  NS-3 body without retraining, dynamic scenario, TCP subflows. The learned pair ranks first of eight.
  Note the inversion of \texttt{single}, mid-table on the mock and last here. Source:
  \texttt{rl-mpquic/runs/eval-ns3/evaluation\_results.json}.}
  \label{tab:p2ns3}
  \begin{tabular}{lccccc}
    \toprule
    Policy & QoE & VMAF & Latency & Loss & Bitrate \\
    \midrule
    Learned pair         & \textbf{0.388} & 79.1 & \SI{84.0}{\milli\second}  & 0.145 & 2702 \\
    Path-only (GCC rate) & 0.340 & 66.3 & \SI{69.9}{\milli\second}  & 0.114 & 1777 \\
    WebRTC               & 0.302 & 56.3 & \SI{57.3}{\milli\second}  & 0.077 & 1424 \\
    App-only             & 0.236 & 77.5 & \SI{100.0}{\milli\second} & 0.247 & 2501 \\
    Proportional         & 0.177 & 50.3 & \SI{62.6}{\milli\second}  & 0.131 & 1485 \\
    Even split           & 0.131 & 25.0 & \SI{41.3}{\milli\second}  & 0.002 & 300 \\
    Random split         & 0.125 & 25.2 & \SI{41.0}{\milli\second}  & 0.003 & 302 \\
    Single best path     & 0.033 & 59.7 & \SI{104.2}{\milli\second} & 0.269 & 1972 \\
    \bottomrule
  \end{tabular}
\end{table}

The ranking transfers: the learned pair is first of eight on a transport stack it never trained against,
and the hierarchy ablation still separates, with app-only at 61\% of the full system. Two differences
from the mock are worth stating plainly rather than burying. Absolute QoE is \emph{not} portable --- the
learned pair scores 0.388 here against 0.617 on the mock, at 0.145 loss against 0.043 --- so the
packet-level body is harsher and cross-body comparisons of absolute QoE are meaningless. And the
\texttt{single} baseline \emph{inverts}: mid-table on the mock, worst of eight here, because one real
TCP subflow cannot absorb the offered rate whereas the mock's unbounded analytical queue lets it. That
inversion is the clearest available evidence that the mock's failure model is optimistic, and it is a
reason to treat the mock as a training surrogate rather than as evidence in its own right.

\gap{this table is a single training seed and five episodes, which is thin for the load it carries ---
it is the chapter's only packet-level evidence. Two runs fix it, in this order. (i) Grade the three
existing dynamic checkpoints on NS-3 (\texttt{RUNBOOK.md} Step~4, evaluation only) so the transfer claim
carries an $\mathit{SD}_{\mathrm{train}}$. (ii) Train natively on NS-3 (\texttt{RUNBOOK.md} Step~5,
measured at $\approx$\,4\,h per seed, the same as the mock, but strictly sequential because the
shared-memory segment name is hard-coded) and compare NS-3-trained against mock-trained, both graded on
NS-3. Either outcome is a result: if the mock-trained policy holds its own, the mock is a validated
cheap surrogate and that is a methodological contribution; if it does not, the sim-to-sim gap is real
and quantified.}

\subsection{Ablations}
\label{sec:p2eval:ablation}

\paragraph{Which half of the hierarchy carries the result?} Both, and in a way that depends on the
intent. Removing the transport agent (app-only) costs 61\% of the QoE under \texttt{interactive} but
only 19\% under \texttt{passive}; removing the application agent (path-only) costs 10\% and 17\%
respectively. Neither half is sufficient, and the joint design's premise ---
\cref{sec:p2design:motivation} --- is what the difference measures.

\paragraph{Domain randomization does not help here.} Randomizing the dynamics hazard rates per training
episode costs 0.027 QoE and loses in all nine evaluation cells, against training-seed spreads of 0.024
and 0.013 --- a real negative result, not noise. The interpretation is mechanical: the observation
carries no signal about \emph{which} dynamics draw is in force, so widening the training distribution
adds variance the policy cannot condition on. Consistent with this, the one place randomization clearly
helps is the app-only fallback (0.299 against 0.155), which is what one expects from a treatment that
flattens a policy rather than sharpening it. One caveat bounds the claim: what is measured is
randomization's \emph{cost} on the nominal environment, because the randomization flag is consumed only
during training and evaluation never resamples. Its intended \emph{benefit} --- generalization to
networks outside the training distribution --- is untested.

\paragraph{An attention-pool critic does not help either, and destabilizes training.} Replacing the
masked mean-pool of \cref{sec:p2design:scoring} with a Set-Transformer attention pool --- targeted at
the shared-bottleneck structure the mean-pool is blind to --- produces a paired difference of 0.083 in
favor of the \emph{mean} pool, inside its own confidence interval and won in only four of nine cells,
i.e.\ not significant. The decisive number is elsewhere: the attention arm's training-seed spread is
0.151, twelve times the mean-pool's 0.013, and larger than the entire margin the system claims over any
baseline. One of its three training seeds collapsed outright, and in all three of that seed's cells the
learned pair is beaten by both WebRTC and the app-only ablation. The mean pool remains the default.
This is worth reporting rather than dropping, because it is a concrete instance of the methodological
point of \cref{sec:p2eval:setup}: evaluated on its healthy seed alone, the attention pool looked like a
small win.

\gap{the flat-versus-scoring comparison is still missing, and it is the ablation that justifies this
chapter's architecture. Train the flat transport agent (\texttt{path\_arch: flat}, fixed path count) on
the dynamic scenario and grade it against the scoring agent under all three intents. The expected
finding is that only the mask-aware scoring agent handles a changing live-path count, since the flat
agent has no representation for a path appearing or disappearing --- but ``expected'' is exactly the
word \cref{ch:single_path} does not have to use, because \cref{tab:p1eval:ablation} measured it. Two
further pieces belong with it, both mirroring evidence \cref{ch:single_path} already carries:
permutation invariance of the split under reordering of the path set (verified in
\texttt{rl-mpquic/tests/test\_scoring\_agent.py} but never reported as an experiment), and QoE against
live-path count, the continuous counterpart of \cref{fig:p1eval:pathcount}. Work item W5 in
\texttt{rl-mpquic/docs/THESIS\_BRIEF.md}.}

\gap{report the per-frame inference cost of the transport agent, which is the binding real-time
constraint of the whole design (\cref{sec:analysis:overhead}) and is currently unstated in both
chapters. Every \texttt{evaluation\_results.json} already records decision time per agent call with
warm-up calls discarded, so this is a read rather than a run. State it against the
\SI{33}{\milli\second} frame budget and against the heuristics' cost, and note that it scales linearly
in the path count because the shared per-path encoder is evaluated once per path.}

\paragraph{DroQ and prioritized replay are unevidenced.} The critic LayerNorm with a raised
update-to-data ratio, and prioritized experience replay, are enabled in every dynamic configuration and
have never been ablated. \gap{run the control arm (\texttt{rl-mpquic/configs/dynamic\_plain.yaml},
which differs in exactly three learner keys, guarded by a test) against the default and report the
paired difference. Note the caveat that shapes how the result should be read: the shipped
update-to-data ratio is 2, not the 10--20 the recipe is about, so a null result would bound a weak
version of the claim. Either report it as such or add a third arm at the higher ratio.}

\subsection{Summary}
\label{sec:p2eval:summary}

Four claims survive this evaluation, and one does not.

The learned pair beats every one of seven comparison points under all three application intents, with
stable rankings and margins of 10.7--20.4\% over the strongest of them, including a WebRTC-style
deployed stack. Its policies behave differently under different intents in the intended direction,
spending a widened delay budget on quality. The ranking transfers zero-shot to a packet-level transport
stack the policies never trained against, though absolute QoE does not. And the scheduler's contribution
scales with both the heterogeneity of path headroom and the strictness of the application's intent ---
the second of which is the axis this chapter shares with \cref{ch:single_path}.

What does not survive is the claim that intents produce intent-specific policies: two of three diagonals
hold, and the third is dominated by a policy trained under a stricter intent, because an objective can
price a term too weakly to learn against it.

Three open items scope all of the above, in decreasing order of what they would change: the chapter has
no intent \emph{conditioning}, only per-intent specialists (\cref{sec:mp:design}); the per-intent
numbers rest on one training seed each; and the packet-level evidence is one seed of one configuration.
